\section{Выводы}
В работе реализован обмен между двумя отдельными программами через отображаемый файл. Родитель управляет вводом и последовательностью запросов, а дочерний процесс проверяет числа и записывает составные в выходной файл. В лабораторной работе №2 данные передавались потоком байтов через каналы, а теперь оба процесса работают с одной областью памяти.

Общая память сама по себе не сообщает, когда данные готовы, поэтому потребовался отдельный механизм уведомления. Пара объектов <<запрос готов>> и <<ответ готов>> превращает обмен в строгое чередование записи: в каждый момент структуру изменяет только один процесс. В Linux для этого использованы семафоры, размещённые прямо в отображённой памяти, а в Windows --- именованные события.

Ожидание без проверки состояния партнёра опасно: гибель одного процесса оставляет второй в бесконечном ожидании. Периодический опрос с проверкой живости процесса решил эту проблему, а тест с уничтожением ребёнка подтвердил, что родитель получает ошибку, а не зависает. Порядок освобождения ресурсов тоже важен: семафор нельзя уничтожать, пока второй процесс может к нему обратиться.

Платформенные различия потребовали аккуратной работы с путями и временем жизни файла. В Windows путь к временной папке может содержать нелатинские символы, а удаление файла по закрытию зависит от режимов разделения. В Linux файл удаляется явно. Разделение на общий интерфейс \texttt{ipc::Channel} и две реализации позволило оставить родительскую и дочернюю программы одинаковыми для обеих систем. Сравнение с лабораторной работой №2 показывает, что отображаемый файл удобен для разделяемых структур фиксированного размера, но накладывает ограничения на размер данных, например на длину строки команды.
